\documentclass[11pt]{letter}
\usepackage[margin=0.9in]{geometry}
\usepackage{amsmath,amssymb}
\usepackage[hidelinks]{hyperref}
\signature{Ramanuja}
\address{Ramanuja\\ Kachhubi, Lakhnaur block\\ Madhubani district, Bihar 847410\\ India\\ \href{mailto:ramanujamath@yahoo.com}{ramanujamath@yahoo.com}}
\date{\today}

\begin{document}
\begin{letter}{The Editor\\ The Ramanujan Journal}

\opening{Dear Editor,}

I am pleased to submit the manuscript entitled \emph{``Proofs of two conjectures of Banerjee and Bringmann on overpartitions and universal mock theta functions''} for consideration for publication in \emph{The Ramanujan Journal}.

In a recent preprint, Banerjee and Bringmann related two families of restricted overpartitions, counted by $\mathrm{op}_1^{[1]}(n)$ and $\mathrm{op}_1^{[2]}(n)$, to Hickerson's universal mock theta functions and conjectured two infinite families of congruences modulo $2$ and $4$. The manuscript proves both conjectures, in each case by establishing a sharper statement: $\mathrm{op}_1^{[1]}(4m)$ is odd exactly when $m$ is triangular, and $\mathrm{op}_1^{[2]}(2m)$ is congruent modulo $4$ to a divisor sum of $8m+1$.

Beyond the conjectures, the paper proves the new congruence $\mathrm{op}_1^{[1]}(4n+1)\equiv0\pmod 4$ and gives an explicit formula for $\mathrm{op}_1^{[1]}(4m)$ modulo $4$ in terms of representation numbers of binary quadratic forms. Combining this with genus theory and a pairing argument in imaginary quadratic orders, it determines the residue modulo $4$ for every prime $p\equiv1\pmod 8$ and every prime $p\equiv13,23\pmod{24}$ exactly dividing $8m+1$, which yields infinitely many explicit congruence families. The methods are elementary $q$-series manipulations together with classical results on mock theta functions and binary quadratic forms, and should be of interest to researchers working on partitions, $q$-series and mock modular forms.

I confirm that the manuscript is original, has not been published previously, and is not under consideration for publication elsewhere. The author declares no competing interests and received no specific funding for this work. As stated in the manuscript, an AI assistant was used for exploratory computations, literature search and drafting; the author has checked all statements and proofs and takes full responsibility for the content.

Thank you for considering this submission. I look forward to hearing from you.

\closing{Sincerely,}

\end{letter}
\end{document}
